% !TEX root = ../../ctfp-print.tex

\lettrine[lhang=0.17]{对}{于范畴里的态射}, 我们可以附上许多种直觉. 但有一点大家都会同意: 如果从物件
$a$ 到物件 $b$ 存在一个态射, 那么这两个物件就以某种方式 ``相关联''. 从某个意义上说, 态射就是这种关
系的证明. 这在任何序集范畴里都一目了然, 那里一个态射 \emph{就是} 一个关系. 一般说来, 两个物件之间的同
一种关系可以有许多 ``证明''. 这些证明组成一个集合, 我们叫它 hom-集合. 当我们变动物件时, 就得到一个从
物件对到 ``证明'' 集合的映射. 这个映射是函子性的 --- 对第一个参数反变, 对第二个参数协变. 我们可以把它
看成是在范畴的物件之间建立了一种全局的关系. 这个关系由 Hom-函子描述:
\[\cat{C}(-, =) \Colon \cat{C}^{op}\times{}\cat{C} \to \Set\]
一般说来, 任何这样的函子都可以解读成在范畴的物件之间建立关系. 一个关系也可能牵涉两个不同的范畴
$\cat{C}$ 和 $\cat{D}$. 描述这种关系的函子有如下签名, 它叫作前函子 (profunctor):
\[p \Colon \cat{D}^{op}\times{}\cat{C} \to \Set\]
数学家说它是从 $\cat{C}$ 到 $\cat{D}$ 的前函子 (注意那个反转), 并且用一个带斜杠的箭头来作它的记号:
\[\cat{C} \nrightarrow \cat{D}\]
你不妨把一个前函子想成 $\cat{C}$ 的物件与 $\cat{D}$ 的物件之间的一种 \newterm{proof-relevant
relation 证明相关的关系}, 其中集合的元素就是这个关系的证明. 每当 $p\ a\ b$ 是空的, $a$ 与 $b$
之间就没有关系. 请记住, 关系未必是对称的.

另一个有用的直觉, 是自函子即容器这一想法的推广. 一个类型为 $p\ a\ b$ 的前函子值, 就可以看成
装满了 $b$ 的容器, 而这些 $b$ 以 $a$ 类型的元素为键. 特别地, hom-前函子的一个元素就是一个从
$a$ 到 $b$ 的函数.

在 Haskell 里, 前函子被定义成一个双参数的类型构造子 \code{p}, 它装备了一个叫 \code{dimap} 的方
法, 后者提升一对函数, 其中第一个函数的方向是 ``反'' 的:

\src{snippet01}
前函子的函子性告诉我们: 如果我们有 \code{a} 与 \code{b} 相关的证明, 那么只要存在一个从 \code{c}
到 \code{a} 的态射, 再加一个从 \code{b} 到 \code{d} 的态射, 我们就得到了 \code{c} 与 \code{d}
相关的证明. 或者, 我们不妨把第一个函数想成是把新的键翻译成旧的键, 把第二个函数想成是在改动容器的内容.

对于只在一个范畴内部起作用的前函子, 我们能从 $p\ a\ a$ 这种对角元素里挖出相当多的信息. 只要有一
对态射 $b \to a$ 和 $a \to c$, 我们就能证明 $b$ 与 $c$ 相关. 更好的是, 我们还能用单个态
射够到非对角位置. 比如, 若有一个态射 $f \Colon a \to b$, 我们可以提升对
$\langle f, \idarrow[b] \rangle$, 从 $p\ b\ b$ 走到 $p\ a\ b$:

\src{snippet02}
或者提升对 $\langle \idarrow[a], f \rangle$, 从 $p\ a\ a$ 走到 $p\ a\ b$:

\src{snippet03}

\section{双自然变换}

既然前函子是函子, 我们就能按标准方式定义它们之间的自然变换. 不过在许多情形里, 只定义两个前函子的
对角元素之间的映射就已经够了. 这样的变换叫作双自然变换, 前提是它满足那些交换条件 —— 这些条件反映的
正是我们把对角元素连到非对角元素上去的两种途径. 两个前函子 $p$ 与 $q$ 之间的一个双自然变换, 其中它
们都属于函子范畴 ${[}\cat{C}^{op}\times{}\cat{C}, \Set{]}$, 就是一族态射:
\[\alpha_a \Colon p\ a\ a \to q\ a\ a\]
使得对任何 $f \Colon a \to b$, 下图都交换:

\begin{figure}[H]
  \centering
  \includegraphics[width=0.35\textwidth]{images/end.jpg}
\end{figure}

\noindent
注意, 这比自然性条件严格地弱. 若 $\alpha$ 是 ${[}\cat{C}^{op}\times{}\cat{C}, \Set{]}$ 里的一个
自然变换, 那么上面这个图就可以由两个自然性方块加上一条函子性条件 (前函子 $q$ 保持复合) 构造出来:

\begin{figure}[H]
  \centering
  \includegraphics[width=0.4\textwidth]{images/end-1.jpg}
\end{figure}

\noindent
注意, ${[}\cat{C}^{op}\times{}\cat{C}, \Set{]}$ 中自然变换 $\alpha$ 的一个分量是以一对物件
$\alpha_{a b}$ 为下标的. 而双自然变换只有一个物件作下标, 因为它只映射各自前函子的对角元素.

\section{端}

现在我们已经准备好, 从范畴论的 ``代数'' 前进到不妨称之为范畴论 ``微积分'' 的部分. 端 (以及余端) 的
微积分从传统微积分里借来了思想, 甚至借来了一些记号. 特别地, 余端可以被理解成一个无穷和, 或者说一个
积分, 而端则类似一个无穷积. 甚至还有一样东西, 它长得像狄拉克 $\delta$ 函数.

端是极限的推广, 只不过把函子换成了前函子. 我们有的不再是锥, 而是楔. 楔的底由前函子 $p$ 的对角元素
构成. 楔的顶点是一个物件 (这里是集合, 因为我们考虑的是 $\Set$ 值前函子), 楔的侧面则是一族函数, 把
顶点映射到底里的那些集合上. 你不妨把这族函数想成一个多态函数 —— 一个在返回类型上多态的函数:
\[\alpha \Colon \forall a\ .\ apex \to p\ a\ a\]
与锥不同, 在楔里面没有任何函数把底的各个顶点连起来. 不过, 正如我们先前见过的, 给定 $\cat{C}$ 里的
任意态射 $f \Colon a \to b$, 我们都能把 $p\ a\ a$ 和 $p\ b\ b$ 连到公共的集合
$p\ a\ b$ 上. 因此我们要求下面这个图交换:

\begin{figure}[H]
  \centering
  \includegraphics[width=0.4\textwidth]{images/end-2.jpg}
\end{figure}

\noindent
这就是所谓的 \newterm{wedge condition 楔条件}. 它可以写成:
\[p\ \idarrow[a]\ f \circ \alpha_a = p\ f\ \idarrow[b] \circ \alpha_b\]
或者用 Haskell 记法:

\src{snippet04}
现在我们可以继续这套普遍构造, 把 $p$ 的端定义为普遍楔 —— 一个集合 $e$, 配上一族函数 $\pi$, 使
得对任何另一个以 $a$ 为顶点、配有族函数 $\alpha$ 的楔, 都存在唯一的函数 $h \Colon a \to e$, 让所
有三角都交换:
\[\pi_a \circ h = \alpha_a\]

\begin{figure}[H]
  \centering
  \includegraphics[width=0.4\textwidth]{images/end-21.jpg}
\end{figure}

\noindent
端的记号是积分号, ``积分变量'' 写在下标的位置:
\[\int_c p\ c\ c\]
$\pi$ 的各个分量叫作端的投影:
\[\pi_a \Colon \int_c p\ c\ c \to p\ a\ a\]
注意, 如果 $\cat{C}$ 是一个离散范畴 (除了恒等态射之外没有别的态射), 那么端就是 $p$ 的所有对角元素
在整个范畴 $\cat{C}$ 上的一个全局的积. 稍后我会向你展示, 在更一般的情形里, 端与这个积之间通过一个
等化子相联系.

在 Haskell 里, 端的公式直接翻译成全称量词:

\src{snippet05}
严格说来, 这只不过是 $p$ 的所有对角元素的积, 但由于
\urlref{https://bartoszmilewski.com/2017-04-11/profunctor-parametricity/}{参数化性}, 楔条件自动满足.
对任何函数 $f \Colon a \to b$, 楔条件读作:

\src{snippet06}
或者, 标上类型注解:

\begin{snipv}
dimap f id\textsubscript{b} . pi\textsubscript{b} = dimap id\textsubscript{a} f . pi\textsubscript{a}
\end{snipv}
等式两边的类型都是:

\src{snippet07}
而 \code{pi} 就是那个多态投影:

\src{snippet08}
这里, 类型推断会自动挑出 \code{e} 的正确分量.

正如我们能够把锥的那一整套交换条件表达成一个自然变换, 同样地, 我们也能把所有楔条件归拢成一个双自然
变换. 为此我们需要把常函子 $\Delta_c$ 推广成常前函子, 它把所有物件对都映到单个物件 $c$, 把所有态
射对都映到这个物件的恒等态射. 一个楔就是从这样一个函子到前函子 $p$ 的双自然变换. 事实上, 一旦我们
意识到 $\Delta_c$ 把所有态射都提升为同一个恒等函数, 双自然性的六边形就收缩成了楔的菱形.

端也可以为 $\Set$ 之外的目标范畴定义, 但在这里我们只考虑 $\Set$ 值前函子以及它们的端.

\section{作为等化子的端}

端的定义里那个交换条件, 可以用一个等化子来写. 首先, 我们来定义两个函数 (这里我用 Haskell 记法, 因
为数学记法在这个例子里似乎更不友好). 这两个函数对应于楔条件中汇聚的那两条分支:

\src{snippet09}[b]
这两个函数都把前函子 \code{p} 的对角元素映射成如下类型的多态函数:

\src{snippet10}
这两个函数的类型不同. 不过, 只要我们造出一个足够大的积类型, 把 \code{p} 的所有对角元素都归拢进去,
就能统一它们的类型:

\src{snippet11}
函数 \code{lambda} 与 \code{rho} 从这个积类型诱导出两个映射:

\src{snippet12}
\code{p} 的端就是这两个函数的等化子. 记住, 等化子挑出的是使两个函数相等的那个最大的子集. 在这里,
它挑出的是所有对角元素的积当中, 使那些楔形图交换的那个子集.

\section{作为端的自然变换}

最重要的端的例子, 大概就是自然变换的集合了. 两个函子 $F$ 与 $G$ 之间的一个自然变换, 是从形如
$\cat{C}(F\ a, G\ a)$ 的 hom-集合里挑出来的一族态射. 若不是自然性条件约束着, 自然变换的集合就会只
是这些 hom-集合的积. 事实上, 在 Haskell 里它正是如此:

\src{snippet13}
它在 Haskell 里行得通, 是因为自然性由参数化性推出. 然而在 Haskell 之外, 并非所有穿过这类 hom-集合
的对角截面都能给出自然变换. 但请注意, 映射
\[\langle a, b \rangle \to \cat{C}(F\ a, G\ b)\]
是一个前函子, 所以研究它的端是有意义的. 这就是楔条件:

\begin{figure}[H]
  \centering
  \includegraphics[width=0.4\textwidth]{images/end1.jpg}
\end{figure}

\noindent
我们就从集合 $\int_c \cat{C}(F\ c, G\ c)$ 里挑一个元素. 两个投影会把这个元素映到某个特定变换的两个
分量上去, 我们就把它们叫作:
\begin{align*}
  \tau_a & \Colon F\ a \to G\ a \\
  \tau_b & \Colon F\ b \to G\ b
\end{align*}
在左边的分支里, 我们用 Hom-函子提升态射对 $\langle \idarrow[a], G\ f \rangle$. 你可能还记得, 这样
的提升是实现为同时做前复合与后复合的. 当它作用在 $\tau_a$ 上时, 被提升的那一对给出:
\[G\ f \circ \tau_a \circ \idarrow[a]\]
图的另一个分支给出:
\[\idarrow[b] \circ \tau_b \circ F\ f\]
二者相等 —— 这是楔条件所要求的 —— 正是 $\tau$ 的自然性条件.

\section{余端}
不出所料, 端的对偶就叫余端. 它构造自一个叫作余楔 (cowedge, 要读成 co-wedge, 而不是 cow-edge) 的
东西的对偶.

\begin{figure}[H]
  \centering
  \includegraphics[width=0.25\textwidth]{images/end-31.jpg}
  \caption{一头带边的奶牛?}
\end{figure}

\noindent
余端的记号是积分号, 只是 ``积分变量'' 写在上标的位置:
\[\int^c p\ c\ c\]
就像端与积相关联一样, 余端与余积相关联, 或者说与和式相关联 (在这方面它确实像一个积分, 毕竟积分就是
和的极限). 我们有的不再是投影, 而入射 (injection) 是从前函子的对角元素出发, 向下走进余端的. 若不是
有那些余楔条件, 我们完全可以说 $p$ 的余端要么是 $p\ a\ a$, 要么是 $p\ b\ b$, 要么是
$p\ c\ c$, 如此等等. 或者我们可以说, 存在某个 $a$, 使得余端就是集合 $p\ a\ a$. 我们在端的定义里
用过的全称量词, 到了余端这里就变成了存在量词.

这就是为什么, 用伪 Haskell 我们会把余端定义成:

\begin{snip}{text}
exists a. p a a
\end{snip}
在 Haskell 里编码存在量词的标准做法, 是使用全称量化的数据构造子. 于是我们可以这样定义:

\src{snippet14}
背后的逻辑是: 无论我们选中哪个 $a$, 都应当能用类型族 $p\ a\ a$ 中任何一个类型的某个值来构造一个
余端.

就像端可以用等化子来定义一样, 余端可以用一个 \newterm{coequalizer 余等化子} 来描述. 所有余楔条件
都可以这样归拢: 对所有可能的函数 $b \to a$, 取一个庞大无比的 $p\ a\ b$ 的余积. 在 Haskell 里,
这会被表达成一个存在类型:

\src{snippet15}
求值这个和式有两条途径, 都是用 \code{dimap} 提升那个函数, 再把它作用到前函子 $p$ 上:

\src{snippet16}
其中 \code{DiagSum} 是 $p$ 的对角元素之和:

\src{snippet17}
这两个函数的余等化子就是余端. 余等化子是这样得到的: 从 \code{DiagSum p} 出发, 把那些由同一个
参数分别经 \code{lambda} 或 \code{rho} 作用而得到的值认同为一. 这里的参数, 是一个由函数
$b \to a$ 与 $p\ a\ b$ 的一个元素组成的对. \code{lambda} 与 \code{rho} 的应用会产生两个可能
不同的 \code{DiagSum p} 的值. 而在余端里, 这两个值被认同为一, 于是余楔条件就自动满足了.

把集合中相关的元素认同为一, 这个过程形式上叫作取商. 要定义一个商, 我们需要一个 \newterm{equivalence
relation 等价关系} $\sim$, 一个自反、对称且传递的关系:
\begin{align*}
   & a \sim a                                                         \\
   & \text{if}\ a \sim b\ \text{then}\ b \sim a                       \\
   & \text{if}\ a \sim b\ \text{and}\ b \sim c\ \text{then}\ a \sim c
\end{align*}
这样的关系把集合切分成一个个等价类. 每个类都由彼此相关的元素组成. 我们从每个类里挑出一个代表元, 就组
成了商集. 一个经典的例子, 就是把有理数定义为整数对, 并配上如下等价关系:
\[(a, b) \sim (c, d)\ \text{iff}\ a * d = b * c\]
很容易验证这是一个等价关系. 一对 $(a, b)$ 被解读为分数 $\frac{a}{b}$, 而分子与分母有公因子的分数
被认同为一. 一个有理数就是这样一些分数构成的等价类.

你可能还记得, 在我们早先讨论极限与余极限时说过, Hom-函子是连续的, 也就是说, 它保持极限. 对偶地, 逆变
的 Hom-函子把余极限变成极限. 这些性质可以推广到端与余端, 它们分别是极限与余极限的推广. 特别地, 我们得
到一个非常有用的恒等式, 它能把余端换算成端:
\[\Set(\int^x p\ x\ x, c) \cong \int_x \Set(p\ x\ x, c)\]
我们到伪 Haskell 里来看一下它:

\begin{snipv}
(exists x. p x x) -> c \ensuremath{\cong} forall x. p x x -> c
\end{snipv}
它告诉我们, 一个接受存在类型的函数, 等价于一个多态函数. 这完全说得通, 因为这样的函数必须准备好处理那
个存在类型里可能编码着的任何一种类型. 同样是这条原理告诉我们, 一个接受和式类型的函数必须实现成 case
语句, 配上一组处理器, 和式里有几种类型就有几个处理器. 在这里, 和式类型换成了余端, 而一族处理器变成了
一个端, 也就是一个多态函数.

\section{忍者米田引理}

米田引理里出现的那个自然变换的集合, 可以用一个端来编码, 于是得到如下表述:
\[\int_z \Set(\cat{C}(a, z), F\ z) \cong F\ a\]
还有一个对偶的公式:
\[\int^z \cat{C}(z, a)\times{}F\ z \cong F\ a\]
这个恒等式让人强烈联想到狄拉克 $\delta$ 函数的公式 (一个函数 $\delta(a - z)$, 或者更准确地说,
是一个广义函数, 它在 $a = z$ 处有一个无穷高的峰). 在这里, Hom-函子扮演的就是 $\delta$ 函数的角色.

这两个恒等式合在一起, 有时被叫作忍者米田引理.

要证明第二个公式, 我们会用到米田嵌入的一个推论, 它说: 两个物件同构, 当且仅当它们的 Hom-函子同构.
换句话说, $a \cong b$ 当且仅当存在一个类型如下的自然变换:
\[[\cat{C}, \Set](\cat{C}(a, -), \cat{C}(b, =))\]
而且它是一个同构.

我们从一个任意的物件 $c$ 出发, 把想证的恒等式的左边放进一个走向 $c$ 的 Hom-函子里:
\[\Set(\int^z \cat{C}(z, a)\times{}F\ z, c)\]
利用连续性的论证, 我们可以把余端换成端:
\[\int_z \Set(\cat{C}(z, a)\times{}F\ z, c)\]
现在我们可以利用积与指数之间的那个伴随:
\[\int_z \Set(\cat{C}(z, a), c^{(F\ z)})\]
我们可以 ``执行积分'', 用米田引理得到:
\[c^{(F\ a)}\]
(注意, 我们用到了米田引理的反变版本, 因为函子 $c^{(F z)}$ 对 $z$ 是反变的.) 这个指数物件同构于
如下的 hom-集合:
\[\Set(F\ a, c)\]
最后, 我们借助米田嵌入得到那个同构:
\[\int^z \cat{C}(z, a)\times{}F\ z \cong F\ a\]

\section{前函子的复合}

让我们进一步探索前函子描述关系这一想法 —— 更准确地说, 是一个证明相关的关系, 也就是说集合
$p\ a\ b$ 代表的是 $a$ 与 $b$ 相关的那些证明组成的集合. 如果我们有两个关系 $p$ 与 $q$, 就能试着
把它们复合起来. 若存在一个中介物件 $c$, 使得 $q\ b\ c$ 与 $p\ c\ a$ 都非空, 我们就说 $a$ 与
$b$ 通过 $q$ 复合 $p$ 而相关. 这个新关系的证明, 就是那些各个单独关系的证明组成的对.
因此, 明白了存在量词对应于一个余端、两个集合的笛卡尔积对应于 ``证明对'' 之后, 我们就能用下面的公式定
义前函子的复合:
\[(q \circ p)\ a\ b = \int^c p\ c\ a\times{}q\ b\ c\]
下面是 \code{Data.Profunctor.Composition} 里等价的 Haskell 定义 (略作改名之后):

\src{snippet18}
这里用的是广义代数数据类型 (generalized algebraic data type), 即 \acronym{GADT} 语法, 在这种语法
里, 一个自由的类型变量 (这里是 \code{c}) 会被自动地存在量化. 于是这个 (未柯里化的) 数据构造子
\code{Procompose} 等价于:

\begin{snip}{text}
exists c. (q a c, p c b)
\end{snip}
这样定义出来的复合, 其单位元就是 Hom-函子 —— 这直接由忍者米田引理推出. 因此, 很自然地要问: 是否存
在这样一个范畴, 其中的态射是前函子? 答案是肯定的, 只是有个前提: 前函子复合的结合律与单位律都只在自然
同构的意义下成立. 这样一个定律只在同构意义下成立的范畴, 叫作双范畴 (bicategory, 它比 $\cat{2}$-范畴
更为一般). 于是我们就有了一个双范畴 $\cat{Prof}$, 其中的物件是范畴, 态射是前函子, 态射之间的态射
(也就是所谓的 2-胞腔) 是自然变换. 实际上我们还能走得更远, 因为除了前函子, 我们还有普通的函子作为范畴
之间的态射. 拥有两种态射的范畴叫作双重范畴 (double category).

前函子在 Haskell 的 lens 库和 arrow 库里都扮演着重要的角色.
